\documentclass[11pt,a4paper]{article}
\usepackage[utf8]{vietnam}
\usepackage{amsmath,amssymb}
\usepackage{tabularx}
\usepackage{booktabs}
\usepackage{multirow}
\usepackage[top=1.0cm, bottom=1.0cm, left=1.2cm, right=1.2cm]{geometry}
\usepackage{fancyhdr}

\pagestyle{fancy}
\fancyhf{}
\rfoot{\small\textit{Trang \thepage}}
\lfoot{\small\textit{Hướng dẫn chấm -- Đề 1 (Toán 7, Giữa HK I)}}
\renewcommand{\headrulewidth}{0pt}

\newcolumntype{Y}{>{\centering\arraybackslash}X}

\begin{document}

\begin{center}
  \textbf{\large ĐÁP ÁN VÀ HƯỚNG DẪN CHẤM -- ĐỀ 1}\\[1mm]
  \textbf{\large KIỂM TRA GIỮA HỌC KÌ I -- MÔN TOÁN 7 -- NĂM HỌC 2026 -- 2027}
\end{center}

\vspace{1mm}
\noindent\textbf{\large PHẦN I. (3,0 điểm) Mỗi câu trả lời đúng được 0,25 điểm.}
\vspace{1mm}

\noindent
\begin{tabularx}{\linewidth}{|c|*{12}{Y|}}
\hline
\textbf{Câu} & 1 & 2 & 3 & 4 & 5 & 6 & 7 & 8 & 9 & 10 & 11 & 12 \\ \hline
\textbf{Đáp án} & \textbf{B} & \textbf{A} & \textbf{D} & \textbf{C} & \textbf{B} & \textbf{A} & \textbf{D} & \textbf{C} & \textbf{B} & \textbf{D} & \textbf{A} & \textbf{C} \\ \hline
\end{tabularx}

\vspace{2.5mm}
\noindent\textbf{\large PHẦN II. (2,0 điểm)}
\vspace{0.5mm}

\noindent\textit{Quy tắc chấm: Đúng 1 ý được 0,1đ; đúng 2 ý được 0,25đ; đúng 3 ý được 0,5đ; đúng 4 ý được 1,0đ.}

\vspace{1mm}
\noindent
\begin{tabularx}{\linewidth}{|c|c|c|X|}
\hline
\textbf{Câu} & \textbf{Ý} & \textbf{Đáp án} & \textbf{Giải thích tóm tắt} \\ \hline
\multirow{4}{*}{\textbf{13}} & a) & \textbf{Sai} & $\sqrt{2}$ là số vô tỉ, không phải số hữu tỉ. \\ \cline{2-4}
& b) & \textbf{Đúng} & $1,(3) = 1,333\dots$ có chu kì là 3. \\ \cline{2-4}
& c) & \textbf{Sai} & $b \approx -1,67; d = 1,5; a \approx 1,414 \Rightarrow b < a < d$. \\ \cline{2-4}
& d) & \textbf{Đúng} & $x - 1/2 = 1,5 \Rightarrow x = 2$ hoặc $x - 1/2 = -1,5 \Rightarrow x = -1$. \\ \hline
\multirow{4}{*}{\textbf{14}} & a) & \textbf{Đúng} & $\widehat{M}_3$ và $\widehat{M}_1$ là hai góc đối đỉnh. \\ \cline{2-4}
& b) & \textbf{Đúng} & Cùng vuông góc với đường thẳng thứ ba thì song song. \\ \cline{2-4}
& c) & \textbf{Sai} & $\widehat{M}_2 = 180^\circ - 50^\circ = 130^\circ \Rightarrow \widehat{N}_2 = 130^\circ$ (đồng vị). \\ \cline{2-4}
& d) & \textbf{Sai} & $\widehat{M}_3 = 50^\circ$; tia phân giác tạo góc $25^\circ \Rightarrow \widehat{MPN} = 25^\circ$ (so le trong). \\ \hline
\end{tabularx}

\vspace{2.5mm}
\noindent\textbf{\large PHẦN III. (2,0 điểm) Mỗi câu đúng được 0,5 điểm.}
\vspace{1mm}

\noindent
\begin{tabularx}{\linewidth}{|c|c|X|}
\hline
\textbf{Câu} & \textbf{Đáp án} & \textbf{Hướng dẫn giải tóm tắt} \\ \hline
\textbf{15} & \textbf{3} & $\sqrt{49} - \sqrt{16} = 7 - 4 = 3$. \\ \hline
\textbf{16} & \textbf{-1} & $\dfrac{3}{5}x = -\dfrac{7}{20} - \dfrac{5}{20} = -\dfrac{3}{5} \Rightarrow x = -1$. \\ \hline
\textbf{17} & \textbf{108} & $a \parallel b \Rightarrow \widehat{F}_1 = \widehat{E}_1 = 72^\circ$ (đồng vị); $\widehat{F}_2 = 180^\circ - 72^\circ = 108^\circ$ (kề bù). \\ \hline
\textbf{18} & \textbf{49} & Cạnh mảnh vườn: $\sqrt{150} \approx 12,247\text{ m}$. Hàng rào: $4\sqrt{150} \approx 48,99 \approx 49\text{ m}$. \\ \hline
\end{tabularx}

\vspace{2.5mm}
\noindent\textbf{\large PHẦN IV. Tự luận (3,0 điểm)}
\vspace{1mm}

\noindent
\begin{tabularx}{\linewidth}{|c|c|X|c|}
\hline
\textbf{Câu} & \textbf{Ý} & \textbf{Nội dung lời giải} & \textbf{Điểm} \\ \hline
\multirow{4}{*}{\textbf{19}} & \multirow{2}{*}{a)} & $\dfrac{5}{7} \cdot \dfrac{3}{11} + \dfrac{5}{7} \cdot \dfrac{8}{11} - \dfrac{12}{7} = \dfrac{5}{7} \cdot \left(\dfrac{3}{11} + \dfrac{8}{11}\right) - \dfrac{12}{7}$ & 0,25 đ \\ \cline{3-4}
& & $= \dfrac{5}{7} \cdot 1 - \dfrac{12}{7} = \dfrac{5 - 12}{7} = -1$. & 0,25 đ \\ \cline{2-4}
& \multirow{2}{*}{b)} & $\left(x - \dfrac{1}{2}\right)^2 = \dfrac{9}{16} \Rightarrow x - \dfrac{1}{2} = \dfrac{3}{4}$ hoặc $x - \dfrac{1}{2} = -\dfrac{3}{4}$. & 0,25 đ \\ \cline{3-4}
& & Với $x - \dfrac{1}{2} = \dfrac{3}{4} \Rightarrow x = \dfrac{5}{4}$; với $x - \dfrac{1}{2} = -\dfrac{3}{4} \Rightarrow x = -\dfrac{1}{4}$. Vậy $x \in \left\{\dfrac{5}{4}; -\dfrac{1}{4}\right\}$. & 0,25 đ \\ \hline

\multirow{4}{*}{\textbf{20}} & \multirow{2}{*}{a)} & Ngày 1 bán: $120 \cdot \dfrac{1}{4} = 30\text{ kg}$, còn lại $120 - 30 = 90\text{ kg}$.\newline Ngày 2 bán: $90 \cdot \dfrac{2}{3} = 60\text{ kg}$. & 0,25 đ \\ \cline{3-4}
& & Sau hai ngày, cửa hàng còn lại: $90 - 60 = 30\text{ (kg)}$. & 0,25 đ \\ \cline{2-4}
& \multirow{2}{*}{b)} & Giá bán mỗi kg ngày thứ ba: $18\,000 \cdot (100\% - 10\%) = 16\,200\text{ (đồng/kg)}$. & 0,25 đ \\ \cline{3-4}
& & Số tiền thu được trong ngày thứ ba: $30 \cdot 16\,200 = 486\,000\text{ (đồng)}$. & 0,25 đ \\ \hline

\multirow{4}{*}{\textbf{21}} & \multirow{2}{*}{a)} & Ta có $\widehat{x'AB} = 180^\circ - \widehat{xAB} = 180^\circ - 120^\circ = 60^\circ$ (hai góc kề bù). & 0,25 đ \\ \cline{3-4}
& & Mà $\widehat{x'AB} = \widehat{ABy} = 60^\circ$, hai góc này ở vị trí so le trong nên $xx' \parallel yy'$. & 0,25 đ \\ \cline{2-4}
& \multirow{2}{*}{b)} & Tia $At$ là phân giác của $\widehat{xAB} \Rightarrow \widehat{tAB} = \dfrac{1}{2}\widehat{xAB} = 60^\circ$.\newline Tia $Bz$ là phân giác của $\widehat{ABy'} \Rightarrow \widehat{ABz} = \dfrac{1}{2}\widehat{ABy'} = \dfrac{1}{2}(180^\circ - 60^\circ) = 60^\circ$. & 0,25 đ \\ \cline{3-4}
& & Do đó $\widehat{tAB} = \widehat{ABz} = 60^\circ$, mà hai góc này ở vị trí so le trong nên $At \parallel Bz$. & 0,25 đ \\ \hline
\end{tabularx}

\end{document}
